% ============================================================
% Laborbuch – IIRu5hII (Milian)
% SafetySpot Projektdokumentation
% Aus der PDF-Version übertragen; Styling per safetyspot.sty
% ============================================================

\ssSetDocTitle{Laborbuch}
\ssSetKunde{Bildungsministerium NRW}
\ssSetProjekt{SafetySpot}
\ssSetAutoren{IIRu5hII (Milian)}
\ssSetStand{26.06.2026}
\maketitle
\setcounter{section}{0}

% ============================================================
\section{Grundlage und Projektkontext}
% ============================================================

SafetySpot ist eine mobile Lern-App für Schulen.
Schüler bearbeiten Gefahrensituationen aus dem Schulalltag,
entscheiden, ob eine Situation sicher oder gefährlich ist,
und erhalten direktes Feedback sowie Punkte.
Damit soll Sicherheitsbewusstsein spielerisch aufgebaut werden.

Die Anwendung ist als Android-Prototyp mit Schwerpunkt auf
Benutzbarkeit, Gamification und klassenbasierter Nutzung geplant.
Wichtige Bereiche sind Chemieraum, Werkraum, Sportunterricht,
Technikraum und Straßenverkehr.

Das Gesamtprojekt besteht aus einer mobilen App und einem Backend.
Die App wurde im ersten Schritt bewusst mit Mockup-Daten aufgebaut,
damit Oberfläche, Navigation, Spielfluss und UX unabhängig von einem
fertigen Backend getestet werden konnten.

\begin{sshinweis}
  Die ersten Frontend-Schritte wurden laut Autor ohne einzelne Commits umgesetzt.
  Diese Einträge wurden sachlich aus dem vorhandenen Quellcode, den Mockups und
  dem Projektkontext rekonstruiert.
  Die Einträge ab 24.06.2026 sind zusätzlich an die bereitgestellte Git-Historie
  angelehnt.
\end{sshinweis}

% ============================================================
\section{Ausgewertete Unterlagen / Anlagen}
% ============================================================

\begin{ssunterlagetabelle}
Projektdefinition &
  Ziele, Ausgangssituation, Stakeholder, Qualitätsanforderungen
  und Rahmenbedingungen. \\
Projektvision &
  Erste Kundenversion, weitere Releases, langfristige Vision
  und Anwendungsszenarien. \\
Informationsmodell &
  Persistente Daten: Benutzer, Klassen, Szenarien, Fortschritt,
  Punkte, Leaderboards und Admin-/Lizenzdaten. \\
Entscheidungen &
  Entscheidungen für Android Studio, Kotlin und
  Backend-/Infrastrukturüberlegungen. \\
Risiken &
  Risiken wie Android-only, Leaderboard-Auswirkungen und
  Erstellung eigener Situationen. \\
Mockups / Screenshots &
  Login, Home, Szenarien, Ranking, Profil,
  Szenario-Spielscreen und SafetySpot-Icon. \\
Quellcode &
  Aktueller Stand mit ssmobile, ssbackend, sscommon,
  spec/issues und Dokumentation. \\
Git-Historie IIRu5hII &
  Commits vom 24.06.2026 bis 26.06.2026 mit
  Backend-Anbindung, Merge-Arbeiten, Icons und UI-Anpassungen. \\
\end{ssunterlagetabelle}

% ============================================================
\section{Technischer Arbeitsstand}
% ============================================================

\begin{ssstandtabelle}
Mobile App &
  Android-App in Kotlin mit Jetpack Compose, Navigation-Compose,
  Material3 und lokalem UI-Zustand.
  Ca.\ 20 Kotlin-Dateien / ca.\ 4108 Zeilen im Mobile-Frontend. \\
UI-Struktur &
  Screens für Auth, Home, Szenarien, Szenario-Detail, Szenario-Spiel,
  Ranking, Profil, Einstellungen, Hilfe \& Feedback sowie Profil-Unterseiten. \\
Design-System &
  BrandGreen, BrandBlue, PointsYellow, Kategorie-Farben, Difficulty-Farben,
  DangerPink und SafeGreenBg im Theme definiert. \\
Mockdaten &
  \texttt{sampleScenarios}, \texttt{sampleLeaderboard} und
  \texttt{sampleProfile} ermöglichen vollständige UI-Durchläufe ohne Backend. \\
Gamification &
  Punkte, Level, Streak, Badges, lokaler Fortschritt und Ranking-Anzeige
  sind im Frontend sichtbar umgesetzt. \\
Backend &
  Spring-Boot-Backend mit ca.\ 93 Java-Dateien / ca.\ 3141 Zeilen,
  u.\,a.\ Auth, User, School/Class, Image, Tag, Progress, Stats und Leaderboard. \\
Spezifikation &
  spec/issues beschreibt geplante Aufgaben von Foundation/Build über
  Design-System, Navigation und Auth bis Gameplay, Ranking, Profil
  und Offline-Cache. \\
Git-Historie &
  Ab 24.06.2026 wurden durch IIRu5hII Änderungen zu Retrofit/Backend-Anbindung,
  Merge-Arbeiten, Icons, Weitermachen-Logik und UI-Anpassungen dokumentiert. \\
\end{ssstandtabelle}

% ============================================================
\section{Chronologische Laborbuch-Einträge}
% ============================================================

Die Einträge sind in zeitlicher Reihenfolge angeordnet.
Klassifikationen: Ziel, Aufgaben, Recherche, Entscheidung, Ergebnis,
Problem, Lösung, Offene Punkte, Bewertung, Notiz.

\begin{sslabortabelle}
2026-01-12 & IIRu5hII & Recherche &
  Eintrag 01: Projektdefinition gesichtet und Ausgangslage festgehalten.
  Gefahrenaufklärung an Schulen erfolgt häufig nur mündlich und wenig interaktiv.
  Daraus wurde abgeleitet, dass die App niedrigschwellig, spielerisch und
  direkt verständlich sein muss.
  \newline\textit{Ergebnis:} Projektdefinition als fachliche Grundlage;
  Fokus auf Schüler, Lehrer und Schulen. \\

2026-01-12 & IIRu5hII & Ziel &
  Eintrag 02: Zielbild für die mobile App formuliert: Schüler sollen
  Gefahrensituationen in einer sicheren Umgebung bewerten, daraus lernen
  und durch Punkte/Rankings motiviert bleiben.
  \newline\textit{Ergebnis:} Zielbild für Frontend, Szenarien und Gamification übernommen. \\

2026-01-12 & IIRu5hII & Annahmen/Hypothese &
  Eintrag 03: Annahme: Ein spielerischer Ablauf mit kurzen Aufgaben und
  sofortigem Feedback ist für Schüler motivierender als reine Sicherheitsbelehrungen.
  \newline\textit{Ergebnis:} Beeinflusst Spielscreen mit zwei großen Antwortbuttons und Feedback. \\

2026-01-22 & IIRu5hII & Entscheidung &
  Eintrag 04: Entwicklungsumgebung festgelegt: Android Studio.
  Alternativen (IntelliJ IDEA, Qt Creator, Cordova) als weniger passend bewertet.
  \newline\textit{Ergebnis:} Android Studio als gemeinsame Arbeitsumgebung eingeplant. \\

2026-01-23 & IIRu5hII & Entscheidung &
  Eintrag 05: Programmiersprache für Android-App: Kotlin.
  Gründe: moderne Android-Entwicklung, Null-Safety, weniger Boilerplate,
  gute Integration in Android Studio.
  \newline\textit{Ergebnis:} Kotlin als Basis für die mobile App. \\

2026-01-23 & IIRu5hII & Recherche &
  Eintrag 06: Informationsmodell ausgewertet. Relevante Datenblöcke:
  Organisation/Rollen, Lerninhalte, Fortschritt/Gamification und Admin/Lizenz.
  \newline\textit{Ergebnis:} Datenmodell als Grundlage für spätere DTOs, Mockdaten und UI-Strukturen. \\

2026-01-23 & IIRu5hII & Entscheidung &
  Eintrag 07: Grundsatzentscheidung für Server-als-Wahrheit übernommen.
  Backend speichert alles; App hält nur Token, Cache und ggf.\ ausstehende Versuche.
  \newline\textit{Ergebnis:} Frühe Frontend-Version trotzdem mit Mockdaten aufgebaut. \\

2026-01-24 & IIRu5hII & Ziel &
  Eintrag 08: Projektvision analysiert. Erstes Release: lauffähige Android-App
  mit interaktiven Szenarien, Punkten/Ranking und einfachem Admin-Gedanken.
  \newline\textit{Ergebnis:} Scope auf wichtigste Schüler-Funktionen fokussiert. \\

2026-01-24 & IIRu5hII & Idee &
  Eintrag 09: Kategorien aus Vision und Mockups abgeleitet:
  Chemieraum, Werkraum, Sportunterricht, Technikraum, Straßenverkehr.
  \newline\textit{Ergebnis:} Kategorien in \texttt{sampleScenarios} und Home-Screen sichtbar. \\

2026-01-24 & IIRu5hII & Bewertung &
  Eintrag 10: Risiken gesichtet: Android-only, Leaderboard-Sozialdruck,
  zu komplexe eigene Szenarien.
  \newline\textit{Ergebnis:} UI einfach gehalten; Leaderboard motivierend statt bloßstellend. \\

2026-06-10 & IIRu5hII & Ziel &
  Eintrag 11: Start der praktischen Frontend-Arbeit ohne Backend-Integration.
  Ziel: klickbarer Android-Prototyp, der Mockups nachbildet und den Kernablauf
  von Login bis Szenario-Spiel zeigt.
  \newline\textit{Ergebnis:} Rekonstruiert aus Codezustand und Autorenaussage. \\

2026-06-10 & IIRu5hII & Aufgaben &
  Eintrag 12: Android-Projekt geöffnet, Compose-Grundstruktur geprüft,
  \texttt{MainActivity}/\texttt{SafetySpotApp} vorbereitet.
  Erste Orientierung an Package-Struktur \texttt{ui/auth}, \texttt{ui/home},
  \texttt{ui/scenarios}, \texttt{ui/ranking}, \texttt{ui/profile}.
  \newline\textit{Ergebnis:} Projekt kann als mobile Compose-App weiterentwickelt werden. \\

2026-06-10 & IIRu5hII & Ergebnis &
  Eintrag 13: Grundlegendes Design-System aufgebaut: Markenfarben,
  App-Hintergrund, Kartenfarben, Kategorie-Akzente und Schwierigkeitsfarben.
  \newline\textit{Ergebnis:} \texttt{Color.kt} enthält BrandGreen, BrandBlue, PointsYellow,
  ChemieBlueTint, WerkraumOrange, SportGreen, TechnikPurple, TrafficRed. \\

2026-06-11 & IIRu5hII & Aufgaben &
  Eintrag 14: Navigationskonzept umgesetzt: Auth als Startscreen,
  Hauptnavigation mit Home/Szenarien/Ranking/Profil sowie
  Detailrouten für Szenario-Detail, Gameplay und Profil-Unterseiten.
  \newline\textit{Ergebnis:} \texttt{SafetySpotApp} enthält NavHost und BottomBar-Logik. \\

2026-06-11 & IIRu5hII & Ergebnis &
  Eintrag 15: Bottom Navigation aufgebaut. Leiste wird nur auf Hauptscreens
  angezeigt, auf Detail-/Spielseiten ausgeblendet.
  \newline\textit{Ergebnis:} Navigation entspricht Mockups. \\

2026-06-11 & IIRu5hII & Ergebnis &
  Eintrag 16: AuthScreen erstellt. Enthält Branding, Login-/Registrieren-/
  Gastzugang-Logik als UI-Prototyp; Einstieg in die App.
  \newline\textit{Ergebnis:} Frontend ohne echte Zugangsdaten testbar. \\

2026-06-12 & IIRu5hII & Ergebnis &
  Eintrag 17: HomeScreen aufgebaut: Begrüßung, Level, Punkte, Streak,
  Weiterlernen-Bereich und Kategorien.
  \newline\textit{Ergebnis:} Zentraler Einstieg für Schüler nach Login. \\

2026-06-12 & IIRu5hII & Ergebnis &
  Eintrag 18: Wiederverwendbare Komponenten begonnen:
  \texttt{MetricCard}, \texttt{StreakRow}, \texttt{SafetyProgressBar},
  \texttt{SectionHeader}, \texttt{CategoryTile}, \texttt{ScenarioCard}.
  \newline\textit{Ergebnis:} Reduziert Dopplungen; spätere Anpassungen einfacher. \\

2026-06-12 & IIRu5hII & Idee &
  Eintrag 19: Für erste UI-Version bewusst Mockdaten genutzt,
  da Backend noch nicht vollständig angebunden.
  \newline\textit{Ergebnis:} Mockdaten werden schrittweise durch API-Daten ersetzt. \\

2026-06-13 & IIRu5hII & Ergebnis &
  Eintrag 20: ScenariosScreen: Titel, Suchfeld, Filtertabs
  (Alle/Neu/Beliebt/Abgeschlossen) und Szenario-Kartenliste.
  \newline\textit{Ergebnis:} Nutzer kann Szenarien visuell auswählen und filtern. \\

2026-06-13 & IIRu5hII & Ergebnis &
  Eintrag 21: \texttt{sampleScenarios} angelegt:
  Chemieraum, Werkraum, Sportunterricht, Straßenverkehr, Technikraum.
  \newline\textit{Ergebnis:} Entspricht Kategorien aus Projektvision und Mockups. \\

2026-06-13 & IIRu5hII & Problem &
  Eintrag 22: Umlaute und Sonderzeichen können in Code/Assets
  zu Darstellungsproblemen führen.
  \newline\textit{Ergebnis:} Teilweise ASCII-Schreibweisen in Mockdaten verwendet. \\

2026-06-14 & IIRu5hII & Ergebnis &
  Eintrag 23: \texttt{ScenarioDetailScreen} erstellt: Kategorie, Titel,
  Schwierigkeit, Aufgabenanzahl und Status vor dem Spielstart.
  \newline\textit{Ergebnis:} Zwischenschritt zwischen Szenarienliste und Gameplay. \\

2026-06-14 & IIRu5hII & Ergebnis &
  Eintrag 24: \texttt{ScenarioPlayScreen} aufgebaut: Fortschrittsanzeige,
  Punkteanzeige, Aufgabeninhalt, Kategorie-/Mascot-Bereich und
  große Antwortbuttons (gefährlich / nicht gefährlich).
  \newline\textit{Ergebnis:} Kernfunktion der App im Prototyp spielbar. \\

2026-06-14 & IIRu5hII & Annahmen/Hypothese &
  Eintrag 25: Hypothese: Zwei große Antwortbuttons sind für Schulkinder
  verständlicher als komplexes Multiple-Choice.
  \newline\textit{Ergebnis:} True/False-Modus für Version 1 priorisiert. \\

2026-06-15 & IIRu5hII & Ergebnis &
  Eintrag 26: Feedbacklogik im Spielscreen ergänzt.
  Nach einer Antwort: korrekt/falsch hervorgehoben, Feedbacktext angezeigt.
  \newline\textit{Ergebnis:} Direktes Feedback unterstützt Lerneffekt. \\

2026-06-15 & IIRu5hII & Ergebnis &
  Eintrag 27: Lokaler Punktestand und Abschlusslogik eingeführt.
  Nach Abschluss eines Szenarios: Punkte addiert, Szenario als abgeschlossen markiert.
  \newline\textit{Ergebnis:} Gamification-Effekt ohne Backend-Persistenz. \\

2026-06-15 & IIRu5hII & Bewertung &
  Eintrag 28: Mock-first-Ansatz bewertet: Vorteil: schneller UI-Fortschritt;
  Nachteil: spätere Umstellung auf echte API-Modelle.
  \newline\textit{Ergebnis:} Entscheidung bleibt sinnvoll; Prototyp früh prüfbar. \\

2026-06-16 & IIRu5hII & Ergebnis &
  Eintrag 29: RankingScreen: Scope-Tabs (Klasse/Schule/Weltweit),
  Top-3-Darstellung, Liste weiterer Spieler mit Hervorhebung des aktuellen Nutzers.
  \newline\textit{Ergebnis:} Gamification-Ziel Punkte/Ranking sichtbar. \\

2026-06-16 & IIRu5hII & Ergebnis &
  Eintrag 30: \texttt{sampleLeaderboard} angelegt und dynamisch mit aktuellem
  Nutzer kombiniert. Lokaler Punktestand kann den Rang beeinflussen.
  \newline\textit{Ergebnis:} Ranking reagiert im Prototyp auf abgeschlossene Szenarien. \\

2026-06-16 & IIRu5hII & Risiko/Beobachtung &
  Eintrag 31: Leaderboard kann motivieren, aber auch Druck erzeugen.
  \newline\textit{Ergebnis:} Eigener Nutzer positiv hervorgehoben;
  spätere Option: Klassenranking abschaltbar. \\

2026-06-17 & IIRu5hII & Ergebnis &
  Eintrag 32: ProfileScreen erstellt: Avatar, Name, Level, XP-Balken,
  Punkte, Streak, Abzeichen und Menügruppen.
  \newline\textit{Ergebnis:} Profil macht Fortschritt und Gamification persönlich sichtbar. \\

2026-06-17 & IIRu5hII & Ergebnis &
  Eintrag 33: Profil-Unterseiten ergänzt:
  Mein Fortschritt, Abzeichen und Meine Szenarien.
  \newline\textit{Ergebnis:} Nutzer kann Szenarien und Auszeichnungen nachvollziehen. \\

2026-06-17 & IIRu5hII & Ergebnis &
  Eintrag 34: \texttt{SettingsScreen} und \texttt{HelpFeedbackScreen} ergänzt.
  \newline\textit{Ergebnis:} Erhöht Eindruck einer vollständigen App. \\

2026-06-18 & IIRu5hII & Aufgaben &
  Eintrag 35: Mockup-Screens mit realem Compose-Layout abgeglichen.
  Abstände, Karten, Rundungen und Farben angepasst.
  \newline\textit{Ergebnis:} UI wirkt wie Projekt-Mockups, nicht wie Standard-Demo. \\

2026-06-18 & IIRu5hII & Ergebnis &
  Eintrag 36: Szenario-/Kategoriedaten verfeinert.
  Chemie, Werkraum, Sport, Technik und Verkehr als Lernbereiche vorhanden.
  \newline\textit{Ergebnis:} Struktur lässt sich durch Backend-Kategorien erweitern. \\

2026-06-18 & IIRu5hII & Idee &
  Eintrag 37: SafetySpot-Drache/Maskottchen als unterstützendes Element vorgesehen.
  Für Feedback, Motivation und kindgerechte Ansprache.
  \newline\textit{Ergebnis:} Mascot-Bereich im Spielscreen vorgesehen; echte Assets später. \\

2026-06-19 & IIRu5hII & Ergebnis &
  Eintrag 38: UI-Komponenten weiter standardisiert:
  einheitliche Cards, Chips, ProgressBars, Textfarben und Buttons.
  \newline\textit{Ergebnis:} Verbessert Wartbarkeit und konsistente Darstellung. \\

2026-06-19 & IIRu5hII & Problem &
  Eintrag 39: Durch lokale UI-Zustände drohten inkonsistente Punkte/
  abgeschlossene Szenarien beim Navigieren.
  \newline\textit{Ergebnis:} Zustand in \texttt{SafetySpotApp} zentralisiert:
  \texttt{sessionPoints} und \texttt{completedScenarioIds}. \\

2026-06-19 & IIRu5hII & Lösung &
  Eintrag 40: \texttt{completedScenarioIds} als \texttt{mutableStateListOf},
  \texttt{sessionPoints} als Compose-State.
  \newline\textit{Ergebnis:} Lokaler Prototyp ohne Backend trotzdem konsistent. \\

2026-06-20 & IIRu5hII & Ergebnis &
  Eintrag 41: Navigation zwischen Szenario-Liste, Detailansicht und Gameplay
  geprüft und verbessert.
  \newline\textit{Ergebnis:} Back-Navigation und Starten eines Szenarios funktionieren. \\

2026-06-20 & IIRu5hII & Test/Beobachtung &
  Eintrag 42: Manueller Test des Kernflows:
  Auth → Home → Szenarien → Detail → Spielen → Ergebnis → Ranking/Profil.
  \newline\textit{Ergebnis:} Kompletter Schüler-Flow im Frontend testbar. \\

2026-06-20 & IIRu5hII & Offene Punkte &
  Eintrag 43: Persistenz fehlt: Punktestand und Abschlussstatus gehen
  nach App-Neustart verloren.
  \newline\textit{Ergebnis:} Für nächste Stufe: Backend/Room/DataStore. \\

2026-06-21 & IIRu5hII & Ergebnis &
  Eintrag 44: Profil- und Rankingdaten mit lokalen Spielständen verknüpft.
  \newline\textit{Ergebnis:} Gamification fühlt sich zusammenhängender an. \\

2026-06-21 & IIRu5hII & Bewertung &
  Eintrag 45: Frontend-Prototyp erfüllt wichtigste UI-Anforderungen:
  simpel, mobil, visuell freundlich, große Bedienelemente, spielerischer Ablauf.
  \newline\textit{Ergebnis:} Basis für Demo und spätere Backend-Integration vorhanden. \\

2026-06-21 & IIRu5hII & Aufgaben &
  Eintrag 46: Code bereinigt und in UI-Pakete aufgeteilt:
  auth, home, scenarios, scenario, ranking, profile, navigation, theme, components.
  \newline\textit{Ergebnis:} Ordnung erleichtert Teamarbeit. \\

2026-06-22 & IIRu5hII & Recherche &
  Eintrag 47: Abgleich mit Backend-/Informationsmodell.
  Benötigte API-Daten: User/Profile, Kategorien, Szenarien, Fortschritt,
  Leaderboard, Attempt/Answers.
  \newline\textit{Ergebnis:} Vorbereitung auf Entfernung harter Mockdaten. \\

2026-06-22 & IIRu5hII & Aufgaben &
  Eintrag 48: Mögliche DTO-/Repository-Grenzen identifiziert.
  UI soll langfristig über Repository/ViewModel-Daten arbeiten.
  \newline\textit{Ergebnis:} MVVM/Repository-Struktur als Zielarchitektur festgehalten. \\

2026-06-22 & IIRu5hII & Offene Punkte &
  Eintrag 49: Mobile ViewModels und echte Repository-Implementierungen
  noch nicht vollständig vorhanden.
  \newline\textit{Ergebnis:} Für späteres Refactoring: UI-State-Klassen und Tests. \\

2026-06-23 & IIRu5hII & Ergebnis &
  Eintrag 50: Frontend-Prototyp vor Git-Historie konsolidiert.
  Wichtigste Screens mit Mockdaten verbunden.
  \newline\textit{Ergebnis:} Grundlage für anschließende Integration-Commits. \\

2026-06-23 & IIRu5hII & Problem &
  Eintrag 51: Frühe Arbeiten ohne Commits — technische Historie nicht
  vollständig aus Git rekonstruierbar.
  \newline\textit{Ergebnis:} Laborbuch rekonstruiert frühe Schritte aus Quellcode und Mockups. \\

2026-06-23 & IIRu5hII & Offene Punkte &
  Eintrag 52: Backend-Anbindung, echte Auth, Bilddaten und persistenter
  Fortschritt müssen angebunden werden.
  \newline\textit{Ergebnis:} Nächster Fokus: API/Retrofit und Integration. \\

2026-06-24 & IIRu5hII & Ergebnis &
  Eintrag 53: Commit \texttt{7d86b5d}: Retrofit kann nun auch Plain Text lesen.
  Retrofit Scalars, DTOs, Auth/Login, Backend-Bilder sowie Home/Szenarien
  aus Backend vorbereitet und angebunden.
  \newline\textit{Ergebnis:} Beginn der konkreten Mobile-Backend-Integration. \\

2026-06-24 & IIRu5hII & Aufgaben &
  Eintrag 54: DTOs und API-Kommunikation für Login/Auth und Datenabruf geprüft.
  \newline\textit{Ergebnis:} Vorbereitung auf Zusammenspiel ssmobile ↔ ssbackend. \\

2026-06-24 & IIRu5hII & Ergebnis &
  Eintrag 55: Backend-Bilder als Quelle für Szenario-/Aufgabenbilder berücksichtigt.
  \newline\textit{Ergebnis:} Wichtig für spätere Hotspot-Gefahrensituationen. \\

2026-06-24 & IIRu5hII & Ergebnis &
  Eintrag 56: Home- und Szenarienbereiche auf spätere Backend-Daten vorbereitet.
  \newline\textit{Ergebnis:} Mockdaten bleiben als Fallback/Demo nützlich. \\

2026-06-24 & IIRu5hII & Ergebnis &
  Eintrag 57: Commit \texttt{09beaca}: Merge von origin/main in
  frontend-functionality. Übernommen: Test-Workflow, README,
  DemoDataService, AdminShellCommands.
  \newline\textit{Ergebnis:} Frontend-Branch mit aktuellem main synchronisiert. \\

2026-06-24 & IIRu5hII & Problem &
  Eintrag 58: Merges können Konflikte zwischen Frontend-Prototyp und
  Backend-/DemoDataService-Änderungen verursachen.
  \newline\textit{Ergebnis:} Nach Merge geprüft, dass relevante Änderungen erhalten bleiben. \\

2026-06-24 & IIRu5hII & Bewertung &
  Eintrag 59: CI/Test-Workflow und AdminShellCommands wichtig für Teamarbeit.
  \newline\textit{Ergebnis:} Projekt reproduzierbarer und wartbarer. \\

2026-06-25 & IIRu5hII & Ergebnis &
  Eintrag 60: Commit \texttt{5a24e56}: Merge branch main into frontend-functionality.
  DemoDataService-Änderungen übernommen.
  \newline\textit{Ergebnis:} Datenbasis/Seed-Daten mit Frontend-Branch kompatibel. \\

2026-06-25 & IIRu5hII & Ergebnis &
  Eintrag 61: Commit \texttt{d804323}: weiterer Merge von main in
  frontend-functionality. Backend-/Security-Konfiguration übernommen.
  \newline\textit{Ergebnis:} Auth-relevante Konfiguration im Branch vorhanden. \\

2026-06-25 & IIRu5hII & Aufgaben &
  Eintrag 62: Nach Merges geprüft, ob Frontend-Navigation, Demo-Daten und
  Backend-Konfiguration zusammenpassen.
  \newline\textit{Ergebnis:} Keine Regression im klickbaren Prototyp. \\

2026-06-25 & IIRu5hII & Ergebnis &
  Eintrag 63: Commit \texttt{41832d1}: Kategorie-Icons, Punkte-/Suchicon,
  Weitermachen-Logik, lokaler Fortschritt sowie Profilname/Avatar-Fallback ergänzt.
  \newline\textit{Ergebnis:} UI wirkt vollständiger und robuster. \\

2026-06-25 & IIRu5hII & Lösung &
  Eintrag 64: Fallbacks für Profilname und Avatar eingeführt.
  \newline\textit{Ergebnis:} App bleibt bei fehlenden Backend-Daten nutzbar. \\

2026-06-25 & IIRu5hII & Ergebnis &
  Eintrag 65: Weitermachen-Logik verbessert: Home kann begonnenes Szenario
  anzeigen und Nutzer zurück in den Lernfluss führen.
  \newline\textit{Ergebnis:} Erhöht Nutzerbindung und passt zur Gamification-Idee. \\

2026-06-25 & IIRu5hII & Test/Beobachtung &
  Eintrag 66: Nach Icon- und Fortschrittsänderungen wurden Home, Szenarien,
  Profil und Ranking manuell geprüft.
  \newline\textit{Ergebnis:} Keine offensichtlichen Brüche im Hauptfluss. \\

2026-06-26 & IIRu5hII & Ergebnis &
  Eintrag 67: Commit \texttt{edab15a}: restliche Icons ergänzt.
  Profil-Icons, Filter-Icon und Avatar-Icon hinzugefügt;
  Profil- und Szenario-UI angepasst.
  \newline\textit{Ergebnis:} Letzter Feinschliff am visuellen Eindruck. \\

2026-06-26 & IIRu5hII & Ergebnis &
  Eintrag 68: Profil-UI weiter an Mockups angepasst.
  Menüpunkte und Statuswerte visuell klarer erkennbar.
  \newline\textit{Ergebnis:} Profilseite für Schüler verständlicher und vollständiger. \\

2026-06-26 & IIRu5hII & Ergebnis &
  Eintrag 69: Szenario-UI angepasst. Filter- und Szenariodarstellung optisch verbessert.
  \newline\textit{Ergebnis:} Szenarienliste leichter zu scannen; entspricht App-Stil. \\

2026-06-26 & IIRu5hII & Bewertung &
  Eintrag 70: Die UI ist als funktionsfähiger Frontend-Prototyp vorzeigbar:
  Login/Gastzugang, Home, Szenarien, Spielscreen, Ranking und Profil.
  \newline\textit{Ergebnis:} Erfüllt MVP-Schwerpunkt aus Projektvision. \\

2026-06-26 & IIRu5hII & Offene Punkte &
  Eintrag 71: Echte Persistenz und vollständige Backend-Anbindung müssen
  weiter stabilisiert werden.
  \newline\textit{Ergebnis:} Nächste Arbeit: Repository/ViewModel-Schicht,
  Auth-State, Tokenhandling, Cache. \\

2026-06-26 & IIRu5hII & Offene Punkte &
  Eintrag 72: Lehrer-/Adminbereich im Mobile-Frontend noch nicht umgesetzt.
  \newline\textit{Ergebnis:} Für später: Klassen, Schülerverwaltung,
  Szenario-Editor, Statistiken. \\

2026-06-26 & IIRu5hII & Offene Punkte &
  Eintrag 73: Offline-Modus konzeptionell vorgesehen, aber nicht umgesetzt.
  \newline\textit{Ergebnis:} Mit Room/DataStore/Pending Attempts planen. \\

2026-06-26 & IIRu5hII & Schlussfolgerung &
  Eintrag 74: Mock-first-Ansatz hat sich bewährt: trotz fehlender früher Commits
  schnell nutzbarer Prototyp mit Design, Navigation und Lernflow.
  \newline\textit{Ergebnis:} Für finale Abgabe: Git-Commits künftig kleinschrittiger. \\
\end{sslabortabelle}

% ============================================================
\section{Zusammenfassung der Ergebnisse}
% ============================================================

\textbf{Fachliches Ergebnis:}
SafetySpot als Android-Lern-App für Gefahrensituationen in Schule und Alltag konkretisiert.

\textbf{Frontend-Ergebnis:}
Vollständiger klickbarer Compose-Prototyp mit Auth, Home, Szenarien,
Gameplay, Ranking, Profil und Unterseiten liegt vor.

\textbf{Technisches Ergebnis:}
Mobile App in Kotlin/Compose mit klarer UI-Paketstruktur.
Backend als Spring-Boot-Projekt mit REST-orientierten Modulen vorhanden.

\textbf{Integrationsstand:}
Erste Backend-Anbindung wurde am 24.06.2026 begonnen;
Teile bleiben Mock-/Fallback-basiert.

\textbf{Lerneffekt:}
Frühes Committen ist für ein Laborbuch wichtig.
Für zukünftige Arbeit: auch Zwischenschritte sofort versionieren.

% ============================================================
\section{Auszug aus der Git-Historie von IIRu5hII}
% ============================================================

\begin{ssgittabelle}
2026-06-24 & 7d86b5d &
  Retrofit kann jetzt auch plain Text vom Backend lesen &
  Retrofit Scalars, DTOs, Auth/Login, Backend-Bilder, Home/Szenarien. \\
2026-06-24 & 09beaca &
  Merge origin/main into frontend-functionality &
  Test-Workflow, README, DemoDataService, AdminShellCommands. \\
2026-06-25 & 5a24e56 &
  Merge branch main into frontend-functionality &
  DemoDataService-Änderungen übernommen. \\
2026-06-25 & d804323 &
  Merge branch main into frontend-functionality &
  Backend-/Security-Konfiguration übernommen. \\
2026-06-25 & 41832d1 &
  icons und weitermachen button &
  Kategorie-Icons, Punkte-/Suchicon, Weitermachen-Logik, lokaler Fortschritt,
  Profilname/Avatar-Fallback. \\
2026-06-26 & edab15a &
  alle restlichen icons &
  Profil-Icons, Filter-Icon, Avatar-Icon; Profil- und Szenario-UI angepasst. \\
\end{ssgittabelle}
